
\documentclass[10pt]{article}
\usepackage[utf8]{inputenc}
\usepackage[T1]{fontenc}
\usepackage{amsmath, amssymb, xcolor, geometry, enumitem, multicol, tcolorbox}

\definecolor{mainblue}{RGB}{0, 102, 204}
\geometry{a4paper, margin=1.2cm, top=1cm, bottom=3cm, footskip=1.2cm}

\begin{document}
\enlargethispage{2.5cm}

% Modern Header
\begin{tcolorbox}[colback=mainblue!5, colframe=mainblue, sharp corners, boxrule=0.5pt]
    \small \textbf{Unit:} Unit 0: Foundations of Algebra \hfill \textbf{Name:} \rule{3cm}{0.4pt} \\
    \small \textbf{Topic:} Variables and Expressions / Order of Operations \hfill \textbf{Date:} \rule{2cm}{0.4pt} \\
    \centering \textbf{\large Challenge (Habanero)}
\end{tcolorbox}

\vspace{0.2cm}

% MCQ Section
\begin{multicols}{2}
\begin{enumerate}[label=\textbf{Q\arabic*.}, leftmargin=*, itemsep=6pt, parsep=0pt]

  \item \small Evaluate the expression: $3 \cdot 2 + 11 - 5$
  \begin{enumerate}[label=(\Alph*), leftmargin=0.6cm, nosep, itemsep=2pt]
    \item 15
    \item 16
    \item 17
    \item 19
    \item 21
  \end{enumerate}

  \item \small A hiker's backpack weighs $8$ pounds. If the hiker adds $2 \cdot 3$ energy bars, each weighing $1$ pound, what is the new weight of the backpack?
  \begin{enumerate}[label=(\Alph*), leftmargin=0.6cm, nosep, itemsep=2pt]
    \item 10
    \item 11
    \item 12
    \item 13
    \item 14
  \end{enumerate}

  \item \small Translate the phrase '5 more than $x$' into an algebraic expression
  \begin{enumerate}[label=(\Alph*), leftmargin=0.6cm, nosep, itemsep=2pt]
    \item $x - 5$
    \item $x + 5$
    \item $5x$
    \item $5 - x$
    \item $x^2 + 5$
  \end{enumerate}

  \item \small Simplify the expression: $2x + 3x - x$
  \begin{enumerate}[label=(\Alph*), leftmargin=0.6cm, nosep, itemsep=2pt]
    \item $3x$
    \item $4x$
    \item $5x$
    \item $6x$
    \item $x^2$
  \end{enumerate}

  \item \small A car's fuel consumption is given by $2 \cdot x + 15$, where $x$ is the number of miles driven. If the car has driven $7$ miles, how much fuel has it consumed?
  \begin{enumerate}[label=(\Alph*), leftmargin=0.6cm, nosep, itemsep=2pt]
    \item 20
    \item 21
    \item 29
    \item 31
    \item 33
  \end{enumerate}

  \item \small Evaluate the expression: $(3 + 2)^2 - 1$
  \begin{enumerate}[label=(\Alph*), leftmargin=0.6cm, nosep, itemsep=2pt]
    \item 15
    \item 16
    \item 24
    \item 25
    \item 30
  \end{enumerate}

  \item \small The population of a species of rabbit is given by $5 \cdot x + 20$, where $x$ is the number of months since the initial population was recorded. If $x = 4$, what is the current population of rabbits?
  \begin{enumerate}[label=(\Alph*), leftmargin=0.6cm, nosep, itemsep=2pt]
    \item 40
    \item 45
    \item 50
    \item 55
    \item 60
  \end{enumerate}

  \item \small Simplify the expression: $9 - 3 + 2 \cdot 4$
  \begin{enumerate}[label=(\Alph*), leftmargin=0.6cm, nosep, itemsep=2pt]
    \item 10
    \item 12
    \item 14
    \item 16
    \item 18
  \end{enumerate}

\end{enumerate}
\end{multicols}

\vspace{-0.2cm}
\hrule
\vspace{0.2cm}
\textbf{\small Open Ended Questions (FRQ)}
\begin{enumerate}[label=\textbf{Q\arabic*.}, start=9, itemsep=5pt, topsep=0pt]

  \item \small A park ranger observes the population of a species of deer and records the data as follows: the population is $25$ deer initially and increases at a rate of $3 \cdot x$ deer per month, where $x$ is the number of months since the initial observation. Write an expression for the population of deer after $x$ months and evaluate it for $x = 6$. Provide your work and final answer. \vspace{2cm}

  \item \small A hiker is on a trail that has an elevation gain of $200$ feet per mile. If the hiker has already hiked $2 \cdot x$ miles and wants to reach the summit, which is $1200$ feet above the starting point, what is the value of $x$? Write an equation representing the situation and solve for $x$. Provide your work and final answer. \vspace{2cm}

\end{enumerate}

\vfill

% Chef's Tips Footer
\begin{tcolorbox}[colback=blue!5, colframe=mainblue!50, title=\small \textbf{Chef's Tips}, fonttitle=\bfseries, bottom=1mm]
\begin{itemize}[noitemsep, topsep=2pt, leftmargin=0.5cm]
  \item \scriptsize Emphasize the importance of following the order of operations when evaluating expressions\item \scriptsize Encourage students to use real-world scenarios to make algebra more relatable and interesting\item \scriptsize Remind students to simplify expressions by combining like terms
\end{itemize}
\end{tcolorbox}

\end{document}
  